\documentclass[12pt]{article}
\usepackage[utf8]{inputenc}

\usepackage[left=2.54cm,right=2.54cm,top=2.54cm,bottom=2.54cm]{geometry}
\usepackage{booktabs}
\usepackage{titlesec}
\usepackage{setspace}
\renewcommand{\arraystretch}{1.15}
\usepackage{amsmath}
\usepackage{adjustbox}
\everymath{\rm}
\usepackage{float}
\usepackage{graphicx}
\usepackage{multirow}
\usepackage{xcolor}

\usepackage[numbers,sort&compress]{natbib}
\bibliographystyle{unsrtnat}
\renewcommand{\refname}{References}


\renewcommand\thesubsection{\arabic{section}.\arabic{subsection}}
\renewcommand\thesubsubsection{\thesubsection.\arabic{subsubsection}}

\titleformat{\section}[block]
  {\normalfont\Large\bfseries\centering}
  {\thesection}
  {1em}
  {}
  [\vspace{2mm}\titlerule]

\usepackage[hidelinks]{hyperref}

\begin{document}

\begin{titlepage}
    \centering
    \vspace*{1in}

    \huge
    \textbf{Developing an autonomous agent for enzyme design}

    \vspace{0.5in}
    \LARGE
    A Proposal for PhD Candidacy

    \vspace{0.5in}
    \Huge
    Ruijie Zhu

    \Large
    \vspace{0.5in}
    \textit{Pritzker School of Molecular Engineering\\ University of Chicago}

    \vspace{1in}

    \Large
    Committee: \\

    \vfill
    April 2026
\end{titlepage}

\newpage
\setlength{\parindent}{0pt}
\setlength{\parskip}{1em}

\section*{Project Summary}
\addcontentsline{toc}{section}{Project Summary}

\subsection*{Overview}
Discovering enzymes for reactions with no natural precedent requires synthesizing heterogeneous information (sequences, structures, mechanisms, substrate chemistry) across inconsistently annotated databases. This project develops an autonomous agent framework with three integrated components: (1) an LLM-driven discovery agent that searches databases and literature, ranking candidates by bond-type compatibility, substrate similarity, and reaction direction; (2) a co-folding pipeline that predicts enzyme--substrate(--cofactor) complexes and extracts binding-site features; (3) experimental validation rounds whose measured activities feed back to refine subsequent candidate selection. From a single target-reaction prompt, the agent autonomously produces a ranked candidate list spanning multiple protein families, evaluated for structural compatibility through co-folding.

\subsection*{Intellectual Merit}
This project advances computational enzyme design in several ways. First, it leverages the reasoning capability of large language models to decompose a target reaction into appropriate database queries, using substrate similarity and bond-type classification rather than sequence homology to identify functionally relevant but evolutionarily distant candidates across databases with diverse coverage and annotation conventions. Second, it develops a structure-aware evaluation pipeline that predicts enzyme--substrate or enzyme--substrate--cofactor complexes and extracts binding geometry and pocket chemistry from the co-folded structures, providing mechanistically grounded structural evidence for assessing candidate compatibility with non-natural substrates. Third, it closes the loop between computational prediction and experimental validation by feeding experimental results to inform future candidate selections.

\subsection*{Broader Impact}
The developed agent framework will lower the barrier of computational enzyme discovery by automating a process that currently requires specialized knowledge of enzyme families, database conventions, and scattered literature across multiple fields. Because the framework requires no retraining or reaction-specific configuration, the same pipeline can be applied to other non-natural transformations across diverse areas of biocatalysis and green chemistry.
\newpage

\section{Introduction}
\label{sec:Introduction}

Enzymes are biological catalysts that accelerate chemical reactions by orders of magnitude through transition-state stabilization. While classical enzymology has emphasized substrate specificity, it is now well established that many enzymes exhibit catalytic promiscuity, the ability to catalyze non-native reactions\cite{khersonsky2010enzyme, copley2003enzymes}. This functional diversity is not incidental: promiscuous activities provide the possibility for functional innovation and offer a powerful foundation for engineering enzymes toward unnatural substrates\cite{hult2007enzyme}.

Identifying which promiscuous enzymes are most likely to accept a given non-natural substrate, however, requires evaluating whether the target substrate is structurally and chemically compatible with the enzyme's binding pocket. Advances in protein structure prediction\cite{jumper2021highly, abramson2024alphafold3, passaro2025boltz} now allow researchers to generate high-confidence structural models across vast regions of sequence space, including enzymes that lack experimental structures, creating an opportunity to computationally evaluate binding pocket properties and catalytic geometry before committing to experimental testing.

A compelling test case for this approach is the enzymatic synthesis of nylons. Among the different types, nylon-6,6 is a fully synthetic polyamide for which no natural biosynthetic pathway exists. It is industrially produced by condensation of adipic acid (a C6 diacid) and hexamethylenediamine (HMDA, a C6 diamine), a reaction that requires harsh conditions\cite{carothers1929studies}. Enzymatic synthesis of nylon-6,6 would enable production under ambient conditions with superior regio- and stereoselectivity. Biology provides a rich repertoire of ATP-dependent amide bond-forming enzymes, including NRPS-independent siderophore (NIS) synthetases\cite{oves2007enzymatic, nusca2012functional} and DdaG-like ligases\cite{hollenhorst2009atp}, that catalyze mechanistically analogous reactions via ATP-dependent carboxylate activation, often acting on structurally diverse carboxylic acids and amine substrates. These enzymes represent promising starting points for engineering amide synthetase activity. While recent work has demonstrated that NIS synthetases such as SfaB can catalyze diad formation between adipic acid and HMDA\cite{qian2025nylon}, catalytic efficiency remains far below what is needed for practical nylon-6,6 synthesis. This underscores the need for a comprehensive search to identify the most promising candidates for subsequent engineering.

A complementary challenge is the enzymatic degradation of nylons, which would enable their sustainable recycling under mild conditions. Conventional nylon recycling relies on harsh acid or base hydrolysis at elevated temperatures, yielding mixed products with limited selectivity\cite{patil2014nylon}. In the case of nylon-6 and nylon-6,6 hydrolysis, biology offers several hydrolase families capable of cleaving aliphatic amide bonds such as Ntn-hydrolase-fold nylon oligomer endohydrolases (NylC, EC 3.5.1.117)\cite{negoro2012nylon} and \(\beta\)-lactamase-fold nylon dimer hydrolases (NylB/NylB\('\), EC 3.5.1.46)\cite{negoro2005three}. NylC is particularly notable as a well characterized enzyme with experimentally confirmed endo-type hydrolysis activity towards both substrates. However, apart from NylC and its homologs\cite{drufva2025identification}, there is very limited exploration of other families of hydrolases that can efficiently catalyze nylon hydrolysis. Therefore, a comprehensive search for the most promising hydrolase candidates provides a parallel test case for the agentic framework to evaluate its ability to identify candidates for mechanistically distinct reaction types.

Recent advances in machine learning and large language models have produced powerful tools for individual stages of the enzyme discovery and design pipeline. Physics-based computational enzyme design, pioneered by theozyme-based active site construction in Rosetta\cite{rothlisberger2008kemp}, enables de novo design of catalytic sites but requires specification of the catalytic geometry and does not address candidate discovery from existing enzyme families. Foundation models for enzyme--reaction matching such as CLEAN\cite{yu2023clean}, CLIPZyme\cite{mikhael2024clipzyme}, and EnzymeCAGE\cite{zheng2026enzymecage} retrieve candidates from large protein databases using learned representations of sequence, structure, and reaction chemistry, but their training data is biased toward well-studied enzyme classes. Generative approaches such as GENzyme\cite{hua2024genzyme} and Riff-Diff\cite{braun2026computational} design de novo enzyme structures conditioned on a target reaction but do not systematically search the existing reservoir of naturally promiscuous enzymes that may already accept non-native substrates. Self-driving laboratory platforms such as SAMPLE\cite{rapp2024selfdriving} and the platform proposed by Singh et al.\cite{singh2025autonomous} automate experimental testing but do not address the upstream problem of identifying which enzyme families to target.

LLM frameworks such as ProteinMCP\cite{xu2026proteinmcp} and Genie-CAT\cite{jacob2025genie} orchestrate specialized protein design tools via natural language but focus on engineering known candidates rather than discovering new ones across protein families. Similarly, LLM-powered literature mining systems such as FuncFetch\cite{smith2025funcfetch} and EnzyExtract\cite{wei2025enzyextract} extract reported enzyme--substrate interactions at scale and cannot identify unstudied enzymes with potential activity on novel substrates. Most recently, Enzyme-tk\cite{mora2026enzymetk} integrates 23 open-source tools into a unified pipeline for enzyme discovery, including ML-based retrieval for target reactions, but relies on pre-defined tool chains rather than adaptive reasoning across heterogeneous databases. What remains missing is a framework that can reason across these heterogeneous sources, decomposing a target reaction into appropriate database queries, interpreting results with inconsistent annotations, and synthesizing evidence spanning database annotations, substrate chemistry, structure, and literature to identify naturally promiscuous enzymes as starting points for engineering toward non-natural reactions.

The absence of such a framework reflects three underlying challenges in non-natural enzyme design. First, candidate enzymes must be identified from families that catalyze mechanistically analogous transformations on structurally different substrates, requiring searches that span multiple EC classes and superfamilies. No single database covers this space adequately: the relevant enzyme families are scattered across resources with inconsistent annotations, and many promising candidates lack EC numbers or exist only as unreviewed TrEMBL entries. Second, evaluating whether a candidate can accept a non-native substrate requires integrating heterogeneous evidence, such as sequence similarity, structural data, catalytic mechanism, substrate binding geometry, and literature on promiscuity. Yet, current computational workflows typically optimize along one of these dimensions in isolation, leading to high false-positive rates when candidates are advanced to experimental testing. Third, computational predictions are rarely validated systematically, and when experimental results contradict predictions, the outcomes are not fed back to refine the computational models. Without this feedback loop, the same prediction errors recur across candidate sets, and the field lacks empirical data on which computational features actually predict activity on non-natural substrates.

These challenges motivate the three aims of this proposal. Aim~1 develops an LLM-driven discovery agent that searches across databases and enzyme families, overcoming the annotation inconsistency and fragmented coverage that limit single-database approaches. Aim~2 builds a structure-aware evaluation pipeline that predicts enzyme--substrate complexes or enzyme--substrate--cofactor complexes via co-folding and extracts mechanistically grounded structural descriptors, including binding quality, pocket chemistry, and catalytic geometry. Aim~3 validates candidates through multiple rounds of experimental testing, feeding activity results back to the agent so it can reason over accumulated positive and negative outcomes to refine candidate selection in subsequent rounds.

\subsection{Enzymatic amide bond synthesis and hydrolysis}
\subsubsection{ATP-dependent amide bond synthesis}
Enzymatic amide bond formation is thermodynamically unfavorable under physiological conditions. Biological systems overcome this barrier by coupling condensation to ATP hydrolysis through two principal activation mechanisms\cite{iyer2009amidoligases, lubberink2023biocatalytic}: adenylation, in which the carboxylate attacks the $\alpha$-phosphate of ATP to form an acyl-AMP intermediate (employed by NIS synthetases)\cite{oves2007enzymatic, yamanaka2008epsilon}, and acyl-phosphate activation, in which $\gamma$-phosphoryl transfer generates an acyl-phosphate intermediate (employed by ATP-grasp fold enzymes)\cite{shomura2012crystal, iyer2009amidoligases}. Both mechanisms activate the carboxylate for nucleophilic attack by the amine, coupling bond formation to the thermodynamically favorable hydrolysis of ATP.

Several enzyme families catalyze mechanistically analogous reactions on substrates related to the nylon-6,6 monomers. The NIS synthetase AsbA condenses citrate with spermidine via adenylation\cite{oves2007enzymatic}, and K311A/M mutations in the related synthetase AsbB increase acceptance of simple linear diamines such as 1,8-diaminooctane and cadaverine relative to the native substrate spermidine\cite{nusca2012functional}. The $\gamma$-glutamylputrescine synthetase PuuA, a glutamine synthetase superfamily member, condenses glutamate with putrescine via ATP-dependent $\gamma$-carboxylate activation\cite{kurihara2008gamma}. Although glutamate and adipic acid differ structurally, the activation mechanism is similar to the target amide bond forming reaction. Notably, PuuA also accepts hexamethylenediamine (HMDA), the diamine monomer of nylon-6,6, as well as cadaverine and spermidine, though with reduced activity relative to the native substrate putrescine\cite{kurihara2008gamma}. These examples of natural substrate promiscuity, particularly PuuA's direct acceptance of HMDA, provide the biochemical foundation for engineering nylon-6,6 amide synthetase activity.

\subsubsection{Amide bond hydrolysis}
Although amide bond hydrolysis is thermodynamically favorable, amide bonds are kinetically stable due to resonance stabilization of the carbonyl, with estimated uncatalyzed half-lives of 350–600 years at neutral pH and 25\textdegree C\cite{radzicka1996rates}. Hydrolytic enzymes overcome this barrier through nucleophilic catalysis: an enzymatic nucleophile attacks the amide carbonyl to form a tetrahedral acyl-enzyme intermediate, which is subsequently hydrolyzed by an activated water molecule.

Two structurally distinct protein folds catalyze aliphatic amide hydrolysis through variations of this mechanism. In the N-terminal nucleophile (Ntn) hydrolase fold, enzyme autocleavage exposes an N-terminal threonine whose side-chain hydroxyl serves as the catalytic nucleophile, activated by the $\alpha$-amino group acting as a general base\cite{negoro2012nylon}. In the $\beta$-lactamase fold (e.g. NylB and NylB'), a catalytic serine activated by a conserved lysine–carboxylate dyad performs the nucleophilic attack\cite{negoro2005three}.

Among known nylon hydrolases, NylC (EC 3.5.1.117) is the best-characterized enzyme with confirmed endo-type activity on both nylon-6 and nylon-6,6 oligomers, cleaving internal amide bonds within the polymer chain\cite{negoro2012nylon}. The nylon dimer hydrolases NylB and NylB' (EC 3.5.1.46) cleave the 6-aminohexanoate linear dimer produced during nylon-6 degradation. These hydrolases provide the enzymatic basis for nylon depolymerization and a mechanistically distinct test case for the proposed framework.


\section{Project objectives}
\label{sec:ProjectObjectives}
The objective of this work is to develop an autonomous agent framework for discovering, structurally modeling, and engineering enzymes for desired chemical reactions, using nylon-6,6 amide synthetases and hydrolases as target systems. To this end, I propose the following three aims:

\textbf{Aim 1: Developing an automated agent for enzyme discovery} \\
Build an LLM-driven agent that autonomously searches sequence, literature, and reaction databases to identify candidate enzymes for a user-specified reaction. The agent decomposes the target reaction chemistry, searches along multiple axes (EC classes, substrate analogues, biosynthetic pathways, protein domain families), and evaluates candidates through a structured evidence framework that assesses bond-type compatibility, substrate similarity, and reaction-direction analysis, producing interpretable, multi-dimensional evidence rather than a single composite score. Structural evaluation from Aim~2 will be integrated into the agent's reasoning loop, and the agent's discovery performance will be benchmarked against data-driven retrieval methods to test whether reasoning-based search outperforms learned representations for non-natural reactions.

\textbf{Aim 2: Structure-aware enzyme candidate evaluation} \\
The second aim builds a structure-aware evaluation framework that assesses catalytic competence from co-folded enzyme-substrate complexes. Candidates identified in Aim~1 are co-folded with target substrates and cofactors (if needed), and the resulting holo structures are evaluated across two complementary dimensions: (1) binding quality, which assesses substrate enclosure, steric plausibility, and substrate positioning inside the pocket; (2) pocket chemistry, which evaluates the compatibility of active-site residue composition with the target substrate. Comparison of apo and holo pocket structures further captures substrate-induced conformational changes that may indicate an induced-fit mechanism. These structural features are integrated with the literature-based evidence from Aim~1 to produce a ranked candidate list for experimental testing.

\textbf{Aim 3: Iterative experimental validation and feedback} \\
The third aim validates candidates from Aims~1--2 through multiple rounds of experimental testing, using activity results to progressively refine the agent's reasoning. In each round, candidates are expressed and assayed for catalytic activity. The experimental outcomes, including both positive and negative results, are fed back to the agent as structured evidence alongside the structural descriptors from Aim~2. The agent synthesizes this accumulated evidence to reason about which structural features correlate with activity, adapting its candidate selection strategy for subsequent rounds rather than applying fixed rules.


\section{Preliminary results}
\label{sec:PreliminaryResults}

The agent framework consists of two integrated pipelines: (1) an LLM-driven discovery agent that searches databases and literature to identify and evaluate candidate enzymes based on function, substrate chemistry and literature evidence; and (2) a structure-aware evaluation pipeline that predicts enzyme--substrate or enzyme--substrate--cofactor complexes via co-folding and extracts binding quality, pocket chemistry, and catalytic geometry. The discovery agent identifies the most promising candidates for structural modeling, the co-folding pipeline then assesses whether the substrates physically fit the active site and whether the pocket environment is chemically compatible with the target substrates.

\subsection{Enzyme discovery agent}
\label{sec:agent_architecture}

\subsubsection{Architecture}
The enzyme discovery agent is implemented as a multi-turn LLM-driven system following the ReAct paradigm\cite{yao2023react}, in which the language model (Claude Sonnet 4.6) alternates between reasoning and tool-calling. The agent has access to specialized tools across four categories. Chemistry tools perform SMILES validation, molecular property calculation, fingerprint similarity via Morgan/ECFP4 fingerprints\cite{rogers2010ecfp}, and physicochemical property calculations including molecular volume, 19 functional group categories, and charge-at-pH estimation. Database search tools query UniProt, KEGG, and InterPro for protein sequences, domain architectures, functional annotations, and reaction metadata (Rhea reaction IDs, KEGG enzyme--reaction mappings). Literature mining tools search PubMed and OpenAlex with semantic re-ranking, citation-network traversal, and named entity extraction. Structured reasoning tools maintain a persistent research notebook and support multi-dimensional candidate evaluation with bond-type classification.

\subsubsection{Enzyme discovery workflow}
\begin{figure}[H]
\centering
\includegraphics[width=\textwidth]{figure/architecture_pro.png}
\caption{Architecture of the enzyme discovery agent. The core control loop (left) orchestrates tool invocation through a registry with search deduplication. Functional tools (center) perform chemistry, database, literature, and reasoning operations, each connecting to external APIs (right). State is persisted in a research notebook across iterations.}
\label{fig:agent_architecture}
\end{figure}

The agent follows a structured discovery workflow. Given a target reaction, the agent (1) decomposes the reaction into its constituent bond types, functional groups, and activation chemistry; (2) enumerates analogous molecules containing the target bond type; (3) searches databases and literature along multiple axes (EC numbers, substrate analogues, bond-type keywords, biosynthetic pathways, gene names, and protein domain families); and (4) evaluates candidates through a structured evidence framework. Several architectural features prevent common failure modes in autonomous search.

\textit{Pre-evaluation gates.} Three mandatory gates ensure thorough coverage before any candidate ranking: (1) automatic named entity extraction from all accumulated paper abstracts, combined with citation-network traversal on the most relevant seed papers, to discover enzyme gene names from related papers that the initial keyword searches missed; (2) batch screening with EC-class triage and bond-type classification to filter mechanistically incompatible candidates before per-candidate evaluation; and (3) homolog expansion for promising singleton protein families, ensuring that InterPro families with only one evaluated member are expanded with additional homologs before final ranking.

\textit{Semantic literature re-ranking.} Literature search results from PubMed and OpenAlex are re-ranked using sentence-transformer embeddings (nomic-embed-text v1.5)\cite{nussbaum2024nomic} to prioritize results most relevant to the target reaction.

\textit{Search deduplication and reflection.} Token-level Jaccard similarity (90\% threshold) prevents redundant database and literature queries across the agent's iterative search process. Every five iterations, a programmatic reflection checkpoint extracts gene names from the agent's research notebook, compares them against the cumulative search history, and injects any unsearched genes back into the agent's prompt as coverage gaps requiring action.

\subsubsection{Candidate evaluation framework}
\label{sec:evaluation_framework}
Once candidates are discovered, they are evaluated using a structured, multi-dimensional evidence framework designed to present interpretable data to the LLM for ranking. For each candidate, the evaluator concurrently gathers data from different sources, including UniProt (EC numbers, catalytic activities, cofactors, mutagenesis variants), InterPro (domain architecture and protein family classification), PubMed (characterization and engineering papers), as well as experimental PDB structures and AlphaFold models when available

Enzyme--substrate compatibility is assessed through three complementary lenses. First, each candidate's EC class and subclass are mapped to the bond type it forms or breaks, its reaction direction, and mechanistic context using a 71-entry lookup table spanning all seven IUBMB enzyme classes, allowing the agent to immediately flag candidates whose native chemistry is incompatible with the target reaction. Second, Tanimoto coefficients between Morgan fingerprints\cite{rogers2010ecfp} (ECFP4) of the target and native substrates provide a baseline structural similarity metric. Third, physicochemical properties (molecular weight, chain length, LogP, topological polar surface area, and functional group composition across 19 categories) are compared between target and native substrates.

For multi-substrate reactions, per-substrate similarities are aggregated using a geometric mean rather than an arithmetic mean, which prevents a strong match to one substrate from masking a poor match to the other. Critical warnings, including bond-type mismatch, one-sided substrate mismatch, severe substrate size mismatch, and reaction-direction incompatibility, automatically place candidates in the lowest ranking tier. Only a specific literature citation demonstrating the enzyme performing the target reaction can override a critical warning.

The LLM synthesizes these evidence dimensions through chemical reasoning, generating two complementary reports: a confidence-weighted ranking that prioritizes experimental evidence (Swiss-Prot annotation, experimental structures, literature depth), and a suitability-focused ranking that prioritizes substrate chemistry compatibility and engineering potential. This design choice is to account for the possibility that the best-characterized enzymes may not necessarily be the most chemically compatible.

\subsection{Agent-driven enzyme discovery}
\label{sec:nylon_discovery}

\subsubsection{Amide synthetase discovery for nylon-6,6}
Given only the high-level input ``Find amide synthetases that ligate hexamethylenediamine and adipic acid,'' the agent autonomously identified ATP-dependent condensation as the relevant activation mechanism and decomposed the problem into four search axes: EC classes for ATP-dependent amide ligases (EC 6.3.2.x, EC 6.3.1.x), systematic enumeration of similar molecules containing aliphatic amide bonds (amino acid derivatives, siderophores, cell wall components, secondary metabolites), protein domain families that catalyze analogous condensations (NIS synthetases, ATP-grasp, Mur ligases), and literature reports of substrate promiscuity in amide-forming enzymes. In the reaction analysis phase, the agent identified five distinct catalytic mechanisms for amide bond formation (acyl-adenylate via NIS synthetases, acyl-phosphate via ATP-grasp enzymes, acyl-CoA thioester, glutamine amidotransferase, and NRPS condensation) and enumerated over 50 biologically synthesized molecules containing aliphatic amide bonds to expand the search space beyond peptide-forming enzymes.

From 1,650 initial candidates, hierarchical filtering through the pre-evaluation gates yielded a final set of 27 candidates from 18 distinct enzyme families. NIS synthetases constituted the largest evaluated group (9 candidates, primarily from the aerobactin IucA/IucC family), followed by ATP-grasp enzymes, Mur ligases, and ANL superfamily adenylation domains. The top five amide synthetase candidates identified by the agent are AsbB, AsbA, PuuA, IucA and IucC.

\autoref{tab:native_substrates} lists the native substrates for all evaluated synthetase candidates. To understand why certain enzyme families are better chemical matches for the target reaction, the native substrate functional group profiles were compared against the target substrate pair (\autoref{fig:substrate_diversity}).

\begin{table}[H]
\centering
\caption{Native substrates of all evaluated synthetase candidates.}
\label{tab:native_substrates}
\begin{adjustbox}{max width=\textwidth}
\small
\begin{tabular}{lllll}
\toprule
\textbf{Enzyme} & \textbf{Accession} & \textbf{Family} & \textbf{Native substrates} & \textbf{Organism} \\
\midrule
\multicolumn{5}{l}{\textit{Target: adipic acid + hexamethylenediamine (HMDA)}} \\
\midrule
AsbA & A0A0F7RJ52 & NIS synthetase & citrate + spermidine & \textit{B. anthracis} \\
AsbB & Q81RQ8 & NIS synthetase & citryl-spermidine + spermidine & \textit{B. anthracis} \\
SbnE & Q2G1M9 & NIS synthetase & citrate + L-2,3-diaminopropanoate & \textit{S. aureus} \\
IucA & Q47316 & NIS synthetase & citrate + $N^6$-acetyl-$N^6$-hydroxylysine & \textit{E. coli} \\
IucC & Q47318 & NIS synthetase & citryl-$N^6$-acetyl-$N^6$-hydroxylysine + $N^6$-acetyl-$N^6$-hydroxylysine & \textit{E. coli} \\
RhbC & Q9Z3R0 & NIS synthetase & --- & \textit{R. meliloti} \\
AlcC & Q7VW17 & NIS synthetase & citrate + cadaverine & \textit{B. pertussis} \\
SfaB & A0A6B5U5X3 & NIS synthetase & --- & \textit{S. aureus} \\
PvsB & Q8GRA8 & NIS synthetase & --- & \textit{V. parahaemolyticus} \\
\midrule
CARNS1 & A5YM72 & ATP-grasp & $\beta$-alanine + L-histidine & \textit{H. sapiens} \\
BacD & P39641 & ATP-grasp & L-alanine + anticapsin & \textit{B. subtilis} \\
DdaF & E2JA31 & ATP-grasp & fumaroyl-diaminopropanoate + L-valine & \textit{E. agglomerans} \\
GshB & P04425 & ATP-grasp & $\gamma$-L-glutamyl-L-cysteine + glycine & \textit{E. coli} \\
DdlB & P07862 & ATP-grasp & D-alanine + D-alanine & \textit{E. coli} \\
\midrule
MurD & P14900 & Mur ligase & UDP-MurNAc-L-Ala + D-glutamate & \textit{E. coli} \\
MurC & P17952 & Mur ligase & UDP-MurNAc + L-alanine & \textit{E. coli} \\
CapB & P96736 & Mur ligase fold & poly-$\gamma$-glutamate primer + L-glutamate & \textit{B. subtilis} \\
\midrule
DdaG & E2JA32 & ANL superfamily & fumarate + L-2,3-diaminopropanoate & \textit{E. agglomerans} \\
IaaL & P18204 & ANL superfamily & indole-3-acetate + L-lysine & \textit{P. savastanoi} \\
Pls & B5BR95 & ANL superfamily & L-lysine + L-lysine & \textit{S. noursei} \\
\midrule
PanC & P31663 & Pantoate ligase & pantoate + $\beta$-alanine & \textit{E. coli} \\
\midrule
PuuA & P78061 & GS superfamily & L-glutamate + putrescine & \textit{E. coli} \\
GshA & P0A6W9 & GS superfamily & L-glutamate + L-cysteine & \textit{E. coli} \\
\midrule
TryS & Q9GT49 & GSP synthetase & glutathione + spermidine & \textit{T. cruzi} \\
GspS & Q63A28 & GSP synthetase & glutathione + spermidine & \textit{B. cereus} \\
\midrule
CphA & P56947/P73833 & Cyanophycin syn. & cyanophycin primer + L-aspartate & \textit{Synechocystis} sp. \\
\midrule
PylC & Q8TUC0 & PylC family & 3-methyl-D-ornithine + L-lysine & \textit{M. acetivorans} \\
\bottomrule
\multicolumn{5}{l}{\footnotesize ---: no catalytic activity annotation in UniProt.} \\
\end{tabular}
\end{adjustbox}
\end{table}

Among the synthetase candidates (\autoref{fig:substrate_diversity}), the target substrates have a minimal functional group profile: only carboxylates, primary amines, and aliphatic carbons. ATP-grasp enzymes such as BacD and DdlB most closely match this profile, while NIS synthetases introduce additional amide and hydroxyl groups, and Mur ligases diverge most strongly with six hydroxyl groups and a phosphate from UDP-sugar chemistry. However, functional group matching captures only one dimension of substrate compatibility. It does not account for substrate size, molecular shape, or the spatial arrangement of functional groups within the binding pocket, limitations that motivate the structure-aware evaluation in Section~\ref{sec:cofolding}.

\begin{figure}[H]
\centering
\includegraphics[width=0.85\textwidth]{figure/substrate_diversity.png}
\caption{Native substrate functional group profiles across synthetase candidates compared to the nylon-6,6 target (adipic acid + HMDA, top row). Each cell shows the count of a given functional group summed across both native substrates. Enzyme families are grouped by brackets.}
\label{fig:substrate_diversity}
\end{figure}


\subsubsection{Amide hydrolase discovery for nylon-6,6}
The agent was also applied to the discovery of enzyme candidates for the complementary reaction: hydrolysis of the nylon-6,6 amide bond, providing a second test case with a mechanistically distinct reaction direction. Given the instruction ``Find amide hydrolases for nylon-6,6 hydrolysis,'' the agent autonomously shifted to hydrolase families (EC~3.5) and searched across amidase-signature-fold enzymes, Ntn-hydrolase-fold endohydrolases, $\beta$-lactamase-fold dimer hydrolases, $\alpha$/$\beta$-hydrolase-fold cutinases, and peptidases. From 933 initial candidates, hierarchical filtering yielded 28 evaluated candidates across 11 enzyme families (\autoref{tab:top_hydrolase}).

The top-ranked candidate was NylC (Q79F77), the only experimentally confirmed nylon-6,6 endo-hydrolase, demonstrating that the agent correctly prioritizes the most functionally relevant enzymes when they exist in the database. The agent distinguished NylC's endo-type cleavage (internal amide bonds) from NylB's exo-type cleavage (terminal dimers only). Beyond the known nylon hydrolases (NylC, NylB, NylB$'$), the remaining candidates come from enzyme families whose native substrates are fundamentally different from nylon-6,6. In nylon-6,6 the repeating amide bonds are uncharged and flanked by aliphatic methylene chains. On the other hand, amidase-signature enzymes natively cleave fatty acid amides (FAAH), N-substituted amides (Aam), or carbamates (AmdA), ceramidases cleave sphingolipid amide bonds flanked by hydroxyl groups, cutinases are primarily esterases with only weak amide activity, and Ntn-hydrolases such as penicillin acylase and bile salt hydrolase act on aromatic or steroid substrates. This pattern contrasts sharply with the synthetase case, where multiple enzyme families naturally process substrates with functional groups similar to adipic acid and HMDA. The scarcity of hydrolases whose native substrates share the uncharged aliphatic character of nylon-6,6 raises the question of whether any of these enzymes could nevertheless accommodate the nylon substrate through pocket compatibility, a question that functional annotations alone cannot answer and that motivates the structural analysis in Section~\ref{sec:cofolding}.

\begin{table}[H]
\centering
\caption{All evaluated hydrolase candidates identified by the discovery agent, grouped by substrate chemistry and literature evidence. Native substrates highlight the substrate mismatch between native and target chemistry. Orthologs of the same enzyme from different organisms are grouped.}
\label{tab:top_hydrolase}
\begin{adjustbox}{max width=\textwidth}
\small
\begin{tabular}{lllll}
\toprule
\textbf{Enzyme} & \textbf{Accession} & \textbf{Family} & \textbf{Native substrate} & \textbf{Organism} \\
\midrule
\multicolumn{5}{l}{\textit{Target: nylon-6,6 amide bond (aliphatic, uncharged, secondary amide)}} \\
\midrule
NylC & Q79F77 & Ntn-hydrolase & nylon-6/6,6 oligomers & \textit{P. ureafaciens} \\
NylC & Q1EPR5 & Ntn-hydrolase & nylon-6/6,6 oligomers & \textit{Agromyces} sp. \\
NylB$'$ & P07062 & $\beta$-lactamase & 6-aminohexanoate dimer & \textit{P. ureafaciens} \\
NylB$'$ & A0A0H5NHB5 & $\beta$-lactamase & 6-aminohexanoate dimer & \textit{N. farcinica} \\
NylB & P07061 & $\beta$-lactamase & 6-aminohexanoate dimer & \textit{P. ureafaciens} \\
NylA & P13398 & Amidase signature & 6-aminohexanoate cyclic dimer & \textit{P. ureafaciens} \\
NylA & P13397 & Amidase signature & 6-aminohexanoate cyclic dimer & \textit{Pseudomonas} sp. \\
\midrule
FAAH & P97612 & Amidase signature & fatty acid amides (C16--C20) & \textit{R. norvegicus} \\
FAAH & O00519 & Amidase signature & fatty acid amides (C16--C20) & \textit{H. sapiens} \\
Aam & K9NBS6 & Amidase signature & N-substituted amides & \textit{R. erythropolis} \\
AmdA & Q9AHE8 & Amidase signature & urethane (carbamate) & \textit{A. tumefaciens} \\
AmdA & P84650 & Amidase signature & short-chain amides & \textit{R. rhodochrous} \\
MalE2 & Q9ZIV5 & Amidase signature & malonamate (C3 diacid amide) & \textit{B. japonicum} \\
\midrule
Ceramidase & Q9I596 & Zn-metalloamidase & sphingolipid amides & \textit{P. aeruginosa} \\
Ceramidase & Q9NR71 & Zn-metalloamidase & sphingolipid amides & \textit{H. sapiens} \\
\midrule
Cut2 & Q6A0I4 & $\alpha$/$\beta$-hydrolase & cutin (polyester) & \textit{T. fusca} \\
Cut1 & E9LVH8 & $\alpha$/$\beta$-hydrolase & cutin (polyester) & \textit{T. cellulosilytica} \\
Cut1 & G8GER6 & $\alpha$/$\beta$-hydrolase & cutin (polyester) & \textit{T. fusca} \\
\midrule
AmiE & P11436 & Nitrilase superfamily & short-chain amides (C2--C3) & \textit{P. aeruginosa} \\
AmiE & Q5Z1U0 & Nitrilase superfamily & short-chain amides & \textit{N. farcinica} \\
AmiE & Q01360 & Nitrilase superfamily & short-chain amides & \textit{R. erythropolis} \\
AmiE & P22984 & Nitrilase superfamily & short-chain amides & \textit{R. erythropolis} \\
Amidase & Q9RQ17 & Nitrilase superfamily & short-chain aliphatic amides & \textit{G. stearothermophilus} \\
Amidase & P95896 & Nitrilase superfamily & short-chain amides & \textit{S. solfataricus} \\
\midrule
HipO & P45493 & M20 peptidase & hippurate (benzoylglycine) & \textit{C. jejuni} \\
BSH & Q06115 & Ntn-hydrolase & conjugated bile acids & \textit{L. plantarum} \\
PGA & P06875 & Ntn-hydrolase & phenylacetyl-amino acids & \textit{E. coli} \\
Subtilisin & P04189 & Subtilase & peptide bonds & \textit{B. subtilis} \\
\bottomrule
\end{tabular}
\end{adjustbox}
\end{table}

\subsection{Enzyme--substrate co-folding}
\label{sec:cofolding}
The discovery agent produces a ranked list of candidate enzymes with accession IDs, evaluation data, and structural availability. This output serves as the input to a structure-aware evaluation pipeline. For the task of finding amide synthetases for nylon-6,6, the 27 candidates from Section~\ref{sec:nylon_discovery} were supplemented with two candidates from prior experimental work: McbA (R4R1U5) and PylC (Q8TUC0), which had been tested for amide diad synthesis on nylon-relevant substrates\cite{qian2025nylon}. After excluding the $\varepsilon$-PL synthetase (B5BR95, 1,319 amino acids) due to its large monomer size, 28 synthetase candidates proceeded to structural evaluation. In parallel, a set of 93 amide hydrolase candidates from experimental nylon degradation studies were also co-folded, along with NylC (Q1EPR4) in its autocleavage form as a positive control benchmark. For each candidate, AlphaFold-predicted or experimentally determined structures were retrieved (if available). All co-folded structures were predicted using Boltz-2\cite{passaro2025boltz}.

\subsubsection{Co-folding pipeline}
\label{sec:cofold_pipeline}
The structural evaluation proceeds through five phases:

\begin{figure}[p]
\centering
\vspace*{\fill}
\includegraphics[height=0.9\textheight]{figure/cofold_pipeline.png}
\vspace*{\fill}
\end{figure}
\begin{figure}[t]
\caption{Co-folding pipeline overview. Phase~0 identifies relevant cofactors and metals via AlphaFill. Phase~1 maps catalytic residues through UniProt annotations and structural homology. Phase~2 detects binding pockets using AlphaFill cofactor transplantation, Foldseek, pyKVFinder, and P2Rank. Phase~3 predicts enzyme--substrate complexes with Boltz-2 using cofactor-guided, pocket-constrained, or unconstrained strategies. Phase~4 scores each complex on iPTM, substrate burial, and pocket chemistry.}
\label{fig:cofold_pipeline}
\end{figure}

\textit{Phase 0: Preparation.} The first ten candidate structures from each enzyme family are queried against AlphaFill\cite{hekkelman2023alphafill} to discover relevant cofactors and catalytic metals. The dominant cofactor group is selected by frequency, with crystallographic variants grouped by physiological form (e.g., ANP, ADP, ATP $\rightarrow$ ATP). Metals are retained when they recurrently co-occur with the dominant cofactor across surveyed candidates. Combined with user-specified substrate SMILES, this produces a per-project configuration that defines the complete set of molecules to include in co-folding.

\textit{Phase 1: Catalytic residue mapping.} For each candidate enzyme, catalytic residues are identified through a combination of UniProt annotation extraction and structural homology. UniProt annotations are classified into core residues (active-site residues) and shell residues (other residues that contact substrates, cofactors, or catalytic metal ions). In parallel, InterPro domain IDs from each candidate are used to query UniProt for reviewed (Swiss-Prot) homologs sharing the same domains, with a fallback to unreviewed (TrEMBL) templates when no reviewed homologs are available. Catalytic site residues from these homologs are mapped onto the candidate's structure through structural alignment using ProDy\cite{bakan2011prody}. This approach is critical for candidates with low sequence identity to characterized enzymes, where sequence-based residue mapping fails. For candidates where direct UniProt active-site annotations are unavailable, a complementary Foldseek\cite{vankempen2024foldseek}-based enhancement step searches the PDB for structural homologs at TM-score $\geq$0.7, transfers active-site residue positions from each matched homolog via structural alignment, and filters biochemically implausible transfers using an LLM that evaluates whether the mapped amino acid identity is consistent with the expected catalytic function of the enzyme family.

\textit{Phase 2: Binding pocket detection.} Binding pocket residues are identified through a two-source cascade. AlphaFill transplants ligands and cofactors from crystallographic homologs into AlphaFold-predicted structures, providing initial cofactor positioning while filtering crystallization artifacts (e.g. glycerol, polyethylene glycol, buffer salts). When AlphaFill lacks a relevant cofactor, Foldseek structural homologs (TM-score $\geq$0.5) provide binding-site annotations via spatial transfer from UniProt. In parallel, pyKVFinder\cite{guerra2021pykvfinder} performs geometric cavity detection with adaptive parameters scaled by protein size (probe radius 8.0~\AA\ for proteins $<$400 residues, 4.0~\AA\ for larger proteins), characterizing pocket boundaries and volumes used for downstream geometry analysis. Candidate pockets are scored by their overlap with the catalytic residues mapped in Phase~1, ensuring that the selected cavity corresponds to the catalytically relevant binding site. When none of the preceding methods produce usable pocket contacts for constraint generation, P2Rank\cite{krivak2018p2rank}, a machine-learning-based pocket predictor, serves as a last-resort fallback.

\textit{Phase 3: Co-folding with Boltz-2.} Enzyme--substrate and enzyme--substrate--cofactor complexes are predicted using Boltz-2 with the AlphaFold-predicted structure provided as a template to fix the protein fold, allowing Boltz-2 to focus computational effort on ligand placement rather than re-folding the protein. For each candidate, the co-folding input specifies: (i) the protein sequence, (ii) the target substrates as SMILES strings, (iii) the ATP cofactor and Mg\textsuperscript{2+} metal ion as CCD codes when applicable, and (iv) pocket constraints derived from the evidence hierarchy established in Phases~1--2. Crystallographic cofactor analogs from AlphaFill are normalized to physiological forms (e.g. ANP/ACP $\rightarrow$ ATP) to model the pre-catalytic state. Catalytic metals detected in AlphaFill contacts are included as separate ligands in the Boltz-2 input, filtered to retain only those whose binding site overlaps the cofactor pocket. 

Three co-folding strategies are employed depending on available evidence: \textit{cofactor-guided} (cofactor constrained to AlphaFill-derived contact residues, substrates dock in proximity), \textit{pocket-constrained} (substrates anchored to residues identified by catalytic site mapping, Foldseek/UniProt binding annotations, or P2Rank pocket prediction), and \textit{unconstrained} (no spatial constraints supplied to the diffusion model, used only when all pocket detection methods fail to produce usable contacts). In all constrained modes, pocket constraints are refined by spatial proximity filtering to ensure all ligands are anchored within a single active site. Five diffusion samples are generated per candidate. For cofactor-dependent enzymes (synthetases), the best model is selected by minimizing substrate--cofactor centroid distances; for cofactor-independent enzymes (hydrolases), the best model is selected by minimizing substrate distance to the predicted catalytic residues identified in Phase~1.

\textit{Phase 4: Quality assessment.} Each co-folded complex is evaluated across two complementary dimensions: binding quality and pocket environment.

Binding quality is assessed by the substrate-specific interface predicted TM-score (iPTM) reported by Boltz-2, which measures protein--substrate binding confidence independently of cofactor contributions. Substrate burial, quantified as the fraction of substrate solvent-accessible surface area (SASA) enclosed by the protein, captures whether the substrate is fully accommodated within the binding pocket or partially exposed at the surface.

Pocket environment is characterized by profiling the amino acid composition of the binding pocket in the co-folded (holo) state, including hydrophobic, charged, and polar residue fractions. For uncharged aliphatic substrates, a high hydrophobic fraction and low charged fraction are expected to favor binding. Comparison between the apo (AlphaFold-predicted) and holo pocket structures provides two additional metrics: the residue Jaccard index (fraction of shared pocket-lining residues between apo and holo structures), which measures whether the substrate occupies the predicted catalytic site or a different surface region, and the fractional cavity volume change, which captures the extent of pocket reorganization upon substrate binding. Together, these metrics distinguish candidates with pre-organized binding sites from those requiring substantial conformational rearrangement.

\subsubsection{Co-folding results}
\label{sec:cofold_results}
The co-folding pipeline was applied to the 28 synthetase candidates discovered by the enzyme discovery agent, and a separate set 93 of hydrolase candidates that were experimentally tested but found inactive towards nylon-6,6 hydrolysis. NylC was also co-folded with the substrate as a benchmark. For co-folding input, the amide synthetase candidates were co-folded with adipic acid, hexamethylenediamine, ATP and Mg\textsuperscript{2+}, and the amide hydrolase candidates were co-folded with the nylon-6,6 trimer. Of the 28 synthetases, 27 used cofactor-guided co-folding (ATP anchored to AlphaFill-derived binding sites with Mg\textsuperscript{2+} coordination, substrates cofolded in proximity); one candidate (iaaL) lacking a cofactor template was co-folded in pocket-constrained mode. \autoref{tab:cofolding_summary} summarizes the co-folding results.

\begin{table}[H]
\centering
\caption{Summary of co-folding results for synthetase and hydrolase candidate sets.}
\label{tab:cofolding_summary}
\begin{adjustbox}{max width=\textwidth}
\begin{tabular}{lcc}
\hline
\textbf{Metric} & \textbf{Synthetases ($n$=28)} & \textbf{Hydrolases ($n$=93)} \\
\hline
Median substrate iPTM & 0.958 & 0.797 \\
Substrate iPTM range & 0.491--0.986 & 0.313--0.924 \\
Median substrate SASA burial (\%) & 89.4 & 62.5 \\
Substrate SASA burial range (\%) & 56.6--98.8 & 20.9--82.5 \\
Candidates with direct UniProt active-/binding-site annotation & 15 & 3 \\
Candidates co-folded with pocket constraints & 28 & 89 \\
\hline
\end{tabular}
\end{adjustbox}
\end{table}

The substrate-specific iPTM, which isolates protein--substrate binding confidence from cofactor contributions where applicable. Although only 3 of the 93 hydrolase candidates had direct UniProt active-site annotations, structural homology transfer via Foldseek provided pocket constraints for 89 of 93 candidates, with only 4 candidates requiring unconstrained co-folding. Because the two candidate sets differ in their main co-folding strategy (cofactor-guided for synthetases vs.\ pocket-constrained for hydrolases), their binding metrics are not directly comparable. Among the synthetases, the high median substrate iPTM (0.958) and SASA burial (89.4\%) indicate that most candidates accommodate the substrates within their predicted catalytic pockets. Among the hydrolases, the lower median substrate iPTM (0.797) and burial (62.5\%) are consistent with pocket chemistries evolved for charged polymeric substrates (e.g., cyanophycin, polyglutamate) that are incompatible with the uncharged aliphatic nylon-6,6 backbone; these low binding metrics persist even among candidates co-folded with active-site constraints, confirming that the substrates mismatch these catalytic pockets.

Among the top candidates by interface quality (\autoref{tab:top_synthetases}), the NIS synthetases (IucA, AsbA, RhbC) and ATP-grasp enzymes (DdlB, Mur ligases) achieve the highest iPTM scores, reflecting the structural compatibility of their cofactor-binding folds with the ATP-dependent substrates. However, iPTM alone is insufficient for predicting catalytic competence on non-native substrates: the discovery agent and co-folding pipeline rank candidates differently, with the discovery agent prioritizing substrate chemistry and literature evidence (favoring AsbB, AsbA, PuuA) while co-folding prioritizes structural fit (favoring IucA, RhbC, AsbA). This divergence motivates the integrated structural evaluation in Aim~2.

\begin{table}[H]
\centering
\caption{Top synthetase candidates ranked by substrate-specific interface confidence (Sub.~iPTM), extracted from the per-chain iPTM matrix to exclude cofactor contributions where applicable. Pocket hydrophobic and charged residue fractions are computed from holo (co-folded) pocket residues.}
\label{tab:top_synthetases}
\begin{adjustbox}{max width=\textwidth}
\begin{tabular}{lllcccc}
\hline
\textbf{Enzyme} & \textbf{Accession} & \textbf{Native substrates} & \textbf{Sub.~iPTM} & \textbf{SASA (\%)} & \textbf{Hydrophobic (\%)} & \textbf{Charged (\%)} \\
\hline
Aerobactin synthase (IucA) & Q47316 & citrate + $N^6$-Ac-$N^6$-OH-lysine & 0.986 & 89.6 & 30.1 & 26.5 \\
Rhizobactin synthase (RhbC) & Q9Z3R0 & --- & 0.985 & 96.9 & 27.1 & 38.6 \\
Spermidine--citrate ligase (AsbA) & A0A0F7RJ52 & citrate + spermidine & 0.983 & 98.7 & 20.7 & 36.2 \\
$\beta$-Carboline amide ligase (McbA) & R4R1U5 & --- & 0.979 & 93.4 & 42.9 & 25.5 \\
L-Dap--citrate ligase (SbnE) & Q2G1M9 & citrate + (S)-2,3-diaminopropanoate & 0.975 & 87.4 & 28.6 & 37.7 \\
D-Ala--D-Ala ligase B (DdlB) & P07862 & D-Ala + D-Ala & 0.974 & 96.5 & 32.6 & 29.2 \\
3-Me-D-Orn--L-Lys ligase (PylC) & Q8TUC0 & (3R)-3-Me-D-Orn + L-Lys & 0.968 & 92.7 & 37.0 & 30.0 \\
UDP-MurNAc--L-Ala ligase (MurC) & P17952 & UDP-MurNAc + L-Ala & 0.966 & 70.5 & 25.0 & 32.7 \\
Dapdiamide A synthase (DdaF) & E2JA31 & fumaroyl-Dap + L-Val & 0.963 & 78.6 & 39.7 & 25.0 \\
UDP-MurNAc-Ala--D-Glu ligase (MurD) & P14900 & UDP-MurNAc-L-Ala + D-Glu & 0.957 & 80.3 & 31.6 & 26.3 \\
\hline
\multicolumn{7}{l}{\footnotesize ---: no catalytic activity annotation in UniProt.} \\
\end{tabular}
\end{adjustbox}
\end{table}

\subsubsection{Pocket chemistry and apo-to-holo comparison}
\label{sec:pocket_chemistry}
To move beyond binding confidence, pocket chemistry was analyzed for both the apo (AlphaFold-predicted) and holo (co-folded) structures for all candidates (\autoref{fig:pocket_chemistry}). 

Pfam-classified synthetase families occupy partially overlapping regions of the pocket-chemistry space (\autoref{fig:pocket_chemistry}A,B). NIS synthetases distinctly cluster at the high-charged end (median $\sim$36\% charged residues vs.\ 23--32\% for other families), reflecting their native binding of charged substrates. At physiological pH, adipic acid is a dicarboxylate anion and HMDA is a diammonium cation, and their charged termini are separated by C4 and C6 methylene units. The ideal pocket must therefore balance cationic residues to coordinate the carboxylate groups of adipate, anionic or polar residues to accommodate the ammonium groups of HMDA without electrostatic repulsion, polar and charged residues to coordinate ATP and Mg\textsuperscript{2+}, and hydrophobic contacts for the aliphatic chains of both substrates. Consistent with this reasoning, SfaB, one of the very few synthetases experimentally demonstrated to catalyze adipic acid-HMDA diad formation\cite{qian2025nylon}, occupies an intermediate position in the pocket-chemistry space ($\sim$33\% charged, $\sim$35\% hydrophobic, $\sim$32\% polar) rather than either extreme. NIS synthetases, whose high density of cationic pocket residues evolved to coordinate the compact tricarboxylate citrate, may create electrostatic repulsion with the protonated ammonium groups of HMDA and their surplus of charged residues would be poorly matched to the aliphatic portions of adipate and HMDA. At the opposite extreme, the most hydrophobic pockets (PanC 48.2\%, PuuA 44.2\%, BacD 40.7\%) may lack sufficient charged residues to coordinate the ionic termini of the substrates and ATP. 

NylC (Q1EPR4), the only experimentally validated nylon hydrolase, was co-folded in its autocleavage form with the precursor split at Thr267 to expose the catalytic nucleophile. NylC is a clear outlier in the pocket-chemistry space (\autoref{fig:pocket_chemistry}C,D): its holo pocket is substantially more hydrophobic (43.8\%) and less charged (12.5\%) than all three hydrolase families (medians 27--33\% charged), consistent with an active site in which charged residues are confined to the catalytic machinery while the remainder of the pocket presents a hydrophobic surface to accommodate the uncharged nylon backbone. The large gap between NylC and the nearest hydrolase family suggests that engineering nylon hydrolase activity from any of these families would require extensive remodeling of the pocket electrostatic environment, not changes to a few contact residues.

To assess whether co-folding places substrates in the expected catalytic pocket, apo (AlphaFold-predicted) and holo (co-folded) structures were compared using the residue Jaccard index (fraction of shared pocket-lining residues). Most candidates show high overlap (e.g., PanC = 0.889, BacD = 0.914), confirming that Boltz-2 consistently places substrates in the AlphaFill-predicted catalytic site with minimal pocket reorganization. However, McbA, despite equally stable pocket geometry (Jaccard = 0.883), is experimentally inactive on nylon substrates\cite{qian2025nylon}, underscoring that correct substrate placement is necessary but not sufficient for catalysis.

Together, these preliminary results establish two capabilities: an autonomous discovery agent that identifies candidates from a single reaction prompt, and a co-folding pipeline that characterizes their binding quality, pocket chemistry, and apo-to-holo stability. They also expose three limitations that motivate the proposed work. First, the discovery agent and structural pipeline operate as independent systems: the agent ranks candidates by substrate chemistry and literature evidence, while co-folding ranks by structural fit, and the two rankings diverge (Aim~1). Second, aggregate descriptors such as iPTM and pocket residue composition are insufficient for predicting catalytic competence: the McbA case demonstrates that correct substrate placement does not guarantee activity. The synthetase pocket-chemistry analysis shows that substrate-level reasoning is needed to interpret what pocket composition is preferred for a given reaction (Aim~2). Third, the number of amide synthetase and hydrolase candidates that demonstrate efficient catalysis towards the target substrates is still low, , and systematic experimental testing is needed both to identify improved candidates and to calibrate which structural descriptors predict activity (Aim~3).

\begin{figure}[H]
\centering
\includegraphics[width=0.85\textwidth]{figure/pocket_chemistry_landscape.png}
\caption{Binding pocket composition (hydrophobic vs.\ charged residue fraction) for co-folded candidates with detectable apo pockets, shown in both apo (AlphaFold-predicted) and holo (co-folded) states. Hydrolase panels include 64 of the 93 candidates; the remaining 29 lacked an evidence-guided apo pocket and are excluded from the chemistry analysis. Enzyme families are colored by Pfam/InterPro classification. (A) Apo pockets of synthetases, (B) holo pockets of synthetases, (C) apo pockets of hydrolases, (D) holo pockets of hydrolases.}
\label{fig:pocket_chemistry}
\end{figure}

\section{Proposed work}
\label{sec:ProposedWork}

\subsection{Aim 1: Developing an automated agent for enzyme discovery}

\textbf{Reaction-network traversal through KBase.} The agent's substrate-driven search currently relies on the LLM to brainstorm natural substrates that contain the target bond type, then searches for enzymes that process those substrates. This brainstorming is incomplete and varies between runs: preliminary runs missed enzymes involved in cyanophycin biosynthesis and siderophore assembly because the LLM failed to generate those substrates. Rather than patching this with additional prompting, the proposed work adds systematic traversal of KBase\cite{arkin2018kbase}, a DOE-funded computational platform that links compounds, reactions, genes, and genomes into a unified graph across over 90,000 annotated microbial genomes\cite{overbeek2014seed}. Given a target reaction class (e.g., amide bond-forming ligases), the agent can traverse KBase's reaction-to-gene mappings to retrieve all gene families annotated as catalyzing reactions in that class, including families the LLM never considered. This deterministic traversal complements the LLM's ability to reason about mechanistic analogy with structured coverage.

\textbf{Sequence-based literature search.} The agent's literature tools start from keywords and follow citation networks, meaning enzymes whose characterization papers use organism-specific gene names unrelated to the target reaction are difficult to reach. PaperBLAST\cite{price2017paperblast}, which links over 400,000 protein sequences to 720,000 scientific articles, provides a complementary sequence-first axis. Once the agent has identified a candidate family, it can query PaperBLAST with a representative sequence to retrieve publications about homologous proteins, surfacing substrate-specificity data and related enzymes in the same fold family that keyword searches would not find.

\textbf{Generalization and benchmarking.} The preliminary results validate the agent on amide synthesis and hydrolysis, two reactions that share the same bond type. To test generalizability, the agent will be applied to 3--5 mechanistically distinct reactions for which the catalyzing enzyme families are well documented in BRENDA and UniProt. The benchmark reactions will be selected to span different catalytic mechanisms and substrate classes, testing whether the agent's search strategy transfers to other reaction types. For each reaction, the reference enzyme families will be established in advance, and recall will be compared against CLIPZyme\cite{mikhael2024clipzyme} (learned reaction-enzyme embeddings) and Enzyme-tk\cite{mora2026enzymetk} (pre-defined tool pipelines). These comparisons establish whether reasoning-based search achieves comparable coverage to learned representations and fixed tool chains on documented reactions, and where it falls short, the failure modes will identify which search axes to strengthen.

\subsection{Aim 2: Structure-aware enzyme candidate evaluation}

The discovery agent and co-folding pipeline rank candidates differently because they evaluate orthogonal evidence (Section~\ref{sec:cofold_results}), and the McbA case shows that correct substrate placement alone does not guarantee activity. The proposed work integrates these two systems and extends the pipeline from candidate evaluation to active-site design.

\textbf{Integrating structural evidence into discovery.} The co-folding pipeline currently produces per-candidate structural metrics that are analyzed post hoc, disconnected from the discovery agent's ranking. To close the loop, the pipeline will return a structured report for each candidate containing binding-quality metrics, per-sub-region pocket-chemistry descriptors (defined below), and catalytic-geometry deviation. The agent treats these structural signals as additional evidence alongside its chemistry and literature scores.

\textbf{Substrate-aware pocket descriptors.} The preliminary pocket-chemistry analysis (Section~\ref{sec:pocket_chemistry}) uses residue composition of the entire pocket, which cannot distinguish whether a residue is positioned near the substrate's reactive groups or elsewhere. To add spatial resolution, the pocket will be decomposed into sub-regions defined by 5\,\AA{} shells around each substrate functional group identified from the co-folded ligand coordinates. Each sub-region is characterized by its residue composition, hydrogen-bond capacity, and local electrostatic character, encoding what pocket chemistry is present at each position along the substrate rather than averaged across the entire cavity. This will provide more fine-grained detail about substrate binding.

\textbf{From structural evaluation to active-site design.} Beyond evaluation, the sub-region decomposition also identifies where a pocket is mismatched to the target substrate, specifying which positions to redesign. LigandMPNN\cite{dauparas2025ligandmpnn} can redesign these positions conditioned on the co-folded substrate placement, preserving catalytic residues while remodeling the surrounding pocket. For example, NIS synthetases have excess charged residues near the substrate mid-chain evolved for citrate binding (Section~\ref{sec:pocket_chemistry}); these could be redesigned while retaining the sub-regions that coordinate ATP and the substrate termini. Designs are verified by re-cofolding with Boltz-2 to confirm substrate positioning is maintained, generating a variant library for Aim~3.

\subsection{Aim 3: Iterative experimental validation and feedback}

Aim~3 closes the loop by experimentally testing candidates and feeding measured activities back to the agent, with the goal of learning what structural and chemical features determine whether an enzyme accepts a non-natural substrate.

\textbf{Empirical calibration of selection criteria.} By experimentally screening a diverse set of candidates and comparing the outcomes with computational modeling results, the experimental results reveal which evidence dimensions are predictive and which are not. Candidates where computation and experiment agree confirm which features are important. Candidates where they disagree are also just as informative. An enzyme that the agent rank highly but that is experimentally inactive points to features the current framework does not capture. Conversely, an enzyme ranked low that shows unexpected activity reveals criteria the agent overweighted or compensatory structural mechanisms it did not consider. Accumulating these cases across diverse enzyme families produces a better understanding of what determines non-natural substrate acceptance, knowledge that currently does not exist in any systematic form and that neither computation nor single-enzyme case studies can provide alone. Critically, this calibration is iterative: the agent uses accumulated evidence to revise its selection criteria and re-rank candidates between rounds, improving candidate recommendations.

\textbf{From case-specific calibration to general principles.} When a single enzyme is found to accept a non-natural substrate, it is unclear whether the features that enabled activity are unique to that protein or would predict activity in other enzymes as well. By testing candidates from multiple structurally distinct protein families on the same target reaction, the experimental data enables a cross-family comparison. This helps delineate whether enzyme activity is due to common structural features or other factors. If common features emerge, they constitute generalizable criteria for predicting non-natural substrate acceptance that the agent can apply to related reaction classes. If the patterns are family-specific, then it indicates that candidate selection needs to be family-aware.

\section{Future Directions}
\label{sec:FutureDirections}
The proposed work develops the framework on two target reactions, but the underlying approach, reasoning-based discovery, structure-aware evaluation, and experimental calibration, is designed to be reaction-agnostic. Applying the calibrated agent to mechanistically distinct reaction classes would test whether the structural insights learned from amide chemistry transfer or whether each reaction class requires its own calibration, a question with broad implications for enzyme engineering. Beyond single reactions, the framework could extend to multi-step pathway design, where the agent identifies or designs an enzyme for each transformation in a synthetic route, integrating retrosynthetic planning with enzyme discovery. As the agent accumulates structure-activity data across reaction classes, the resulting knowledge base would capture empirical relationships between active-site structure and non-natural substrate acceptance at a scale that individual case studies cannot achieve, deepening our understanding of the structural basis of enzyme promiscuity.

Looking further ahead, the iterative cycle of computational prediction and experimental testing developed in this work is a step toward fully autonomous enzyme discovery. The integration with automated experimental platforms would further allow the agent to design, execute, and learn from experiments with minimal human intervention. 
\clearpage
\bibliography{citation}

\end{document}
